\section{Representation by $\delta[n]$ and $\mu[n]$; more convolution (2.43--2.49)}

\Prob{2.43\quad The Problem-2.3 sequences as combinations of $\delta$'s and $\mu$'s}

\textbf{(a) Delayed unit samples.} Any finite sequence is
$s[n]=\sum_k s[k]\,\delta[n-k]$, so one simply reads off the sample values:
\[
\boxed{
\begin{aligned}
x[n]&=3\delta[n+1]-4\delta[n]-2\delta[n-1]+5\delta[n-2]+\delta[n-4]-3\delta[n-5],\\
y[n]&=-5\delta[n+3]-\delta[n+2]+3\delta[n]+6\delta[n-1]-2\delta[n-2]+8\delta[n-3],\\
w[n]&=-\delta[n-3]+3\delta[n-4]-4\delta[n-5]+\delta[n-7]-5\delta[n-8]+6\delta[n-9].
\end{aligned}}
\]
(The zero-valued samples $x[3]=0$, $y[-1]=0$, $w[6]=0$ contribute no term.)

\textbf{(b) Delayed unit steps.} Since $\delta[n-k]=\mu[n-k]-\mu[n-k-1]$, substituting into
part (a) and collecting terms telescopes to
\[
s[n]=\sum_{k}\bigl(s[k]-s[k-1]\bigr)\,\mu[n-k],
\]
i.e.\ the step coefficients are the \emph{first differences} of the sample values (one extra
term is needed at the far end to bring the sequence back to zero). Computing the differences:
\[
\boxed{\small
\begin{aligned}
x[n]&=3\mu[n+1]-7\mu[n]+2\mu[n-1]+7\mu[n-2]-5\mu[n-3]+\mu[n-4]-4\mu[n-5]+3\mu[n-6],\\[2pt]
y[n]&=-5\mu[n+3]+4\mu[n+2]+\mu[n+1]+3\mu[n]+3\mu[n-1]-8\mu[n-2]+10\mu[n-3]-8\mu[n-4],\\[2pt]
w[n]&=-\mu[n-3]+4\mu[n-4]-7\mu[n-5]+4\mu[n-6]+\mu[n-7]-6\mu[n-8]+11\mu[n-9]-6\mu[n-10].
\end{aligned}}
\]
Check for $x$: the coefficients are
$3,-7,2,7,-5,1,-4,3$ at $k=-1,\dots,6$; their running sums reproduce
$3,-4,-2,5,0,1,-3,0$, which is $x[n]$ followed by zero. \checkmark
Each coefficient list sums to zero, as it must for a finite-length sequence.

\Prob{2.44\quad $x[n]=\{2.1,\,-3.3,\,-1.7,\,5.2\}$, $-1\le n\le 2$, in terms of $\mu[n]$}

Using the first-difference rule of Problem 2.43(b) with
$x[-1]=2.1$, $x[0]=-3.3$, $x[1]=-1.7$, $x[2]=5.2$ (and $x=0$ outside):
\[
\begin{array}{r|rrrrr}
k & -1 & 0 & 1 & 2 & 3\\\hline
x[k]-x[k-1] & 2.1 & -5.4 & 1.6 & 6.9 & -5.2
\end{array}
\]
\[
\boxed{\;x[n]=2.1\,\mu[n+1]-5.4\,\mu[n]+1.6\,\mu[n-1]+6.9\,\mu[n-2]-5.2\,\mu[n-3].\;}
\]
Check at $n=1$: $2.1-5.4+1.6=-1.7$ \checkmark; at $n\ge3$: $2.1-5.4+1.6+6.9-5.2=0$ \checkmark.

\Prob{2.45\quad Express $x[n]=\begin{cases}2,& n\ge2\\ 1,& n<2\end{cases}$ in terms of $\mu[n]$}

The sequence equals a constant $1$ for all $n$, plus an extra $1$ switched on at $n=2$.
Since the constant sequence $1$ can be written as $\mu[n]+\mu[-n-1]$, or simply carried as a
constant,
\[
\boxed{\;x[n]=1+\mu[n-2].\;}
\]
Verification: for $n<2$, $\mu[n-2]=0$ so $x[n]=1$; for $n\ge2$, $\mu[n-2]=1$ so $x[n]=2$.
\checkmark
(Note that $x[n]$ is not a finite-length or right-sided sequence, so it cannot be written as a
finite combination of \emph{delayed} steps alone --- the additive constant is essential.)

\Prob{2.46\quad Closed-form convolution sums}

\textbf{(a)} $y[n]=\alpha^{n}\mu[n]\circledast\mu[n]$. For $n\ge0$,
\[
y[n]=\sum_{k=-\infty}^{\infty}\alpha^{k}\mu[k]\,\mu[n-k]=\sum_{k=0}^{n}\alpha^{k},
\]
because the two steps restrict $k$ to $0\le k\le n$; for $n<0$ the ranges do not overlap and
$y[n]=0$. Summing the geometric series,
\[
\boxed{\;y[n]=\alpha^{n}\mu[n]\circledast\mu[n]=\frac{1-\alpha^{\,n+1}}{1-\alpha}\,\mu[n]
\qquad(\alpha\ne1),\qquad y[n]=(n+1)\mu[n]\ \ (\alpha=1).\;}
\]
For $|\alpha|<1$ this tends to the steady-state value $1/(1-\alpha)$ --- the DC gain of a
one-pole filter.

\textbf{(b)} $y[n]=n\alpha^{n}\mu[n]\circledast\mu[n]=\displaystyle\sum_{k=0}^{n}k\alpha^{k}$
for $n\ge0$. Differentiating the geometric sum
$\sum_{k=0}^{n}\alpha^{k}=\frac{1-\alpha^{n+1}}{1-\alpha}$ and multiplying by $\alpha$
(i.e.\ applying $\alpha\frac{d}{d\alpha}$),
\[
\sum_{k=0}^{n}k\alpha^{k}
=\alpha\,\frac{d}{d\alpha}\!\left(\frac{1-\alpha^{n+1}}{1-\alpha}\right)
=\frac{\alpha\bigl[\,1-(n+1)\alpha^{n}+n\alpha^{\,n+1}\bigr]}{(1-\alpha)^{2}} .
\]
\[
\boxed{\;y[n]=n\alpha^{n}\mu[n]\circledast\mu[n]
=\frac{\alpha\bigl[\,1-(n+1)\alpha^{n}+n\alpha^{\,n+1}\bigr]}{(1-\alpha)^{2}}\;\mu[n]
\qquad(\alpha\ne1),\;}
\]
and for $\alpha=1$, $y[n]=\sum_{k=0}^{n}k=\dfrac{n(n+1)}{2}\mu[n]$.
For $|\alpha|<1$ the limit is $\alpha/(1-\alpha)^{2}$.

\Prob{2.47\quad Convolutions of four impulse-train sequences}

First list the sequences as sample values (the arrow marks $n=0$):
\[
\begin{aligned}
x_1[n]&=2\delta[n-1]-2\delta[n+1] &&\Rightarrow&& \{-2,\;\underset{\uparrow}{0},\;2\},&& -1\le n\le 1,\\
x_2[n]&=3\delta[n-2]-\delta[n]   &&\Rightarrow&& \{\underset{\uparrow}{-1},\;0,\;3\},&& \ \ 0\le n\le 2,\\
h_1[n]&=\delta[n+3]-1.5\delta[n]-\delta[n-2] &&\Rightarrow&& \{1,0,0,\underset{\uparrow}{-1.5},0,-1\},&& -3\le n\le 2,\\
h_2[n]&=-\delta[n+1]+2\delta[n-1]+3\delta[n-3] &&\Rightarrow&& \{-1,\underset{\uparrow}{0},2,0,3\},&& -1\le n\le 3 .
\end{aligned}
\]
Convolution of impulse trains is easiest term-by-term, using
$a\delta[n-p]\circledast b\delta[n-q]=ab\,\delta[n-p-q]$.

\textbf{(a)} $y_1[n]=x_1[n]\circledast h_1[n]$. The four cross products are
\[
\begin{aligned}
&(2\delta[n-1])\circledast(\delta[n+3])=2\delta[n+2], &&(2\delta[n-1])\circledast(-1.5\delta[n])=-3\delta[n-1],\\
&(2\delta[n-1])\circledast(-\delta[n-2])=-2\delta[n-3], &&(-2\delta[n+1])\circledast(\delta[n+3])=-2\delta[n+4],\\
&(-2\delta[n+1])\circledast(-1.5\delta[n])=3\delta[n+1], &&(-2\delta[n+1])\circledast(-\delta[n-2])=2\delta[n-1].
\end{aligned}
\]
Collecting ($-3\delta[n-1]+2\delta[n-1]=-\delta[n-1]$):
\[
\boxed{\;y_1[n]=-2\delta[n+4]+2\delta[n+2]+3\delta[n+1]-\delta[n-1]-2\delta[n-3],\;}
\]
\[
\{y_1[n]\}=\{-2,\;0,\;2,\;3,\;\underset{\uparrow}{0},\;-1,\;0,\;-2\},\qquad -4\le n\le 3 .
\]

\textbf{(b)} $y_2[n]=x_2[n]\circledast h_2[n]$:
\[
\boxed{\;y_2[n]=\delta[n+1]-5\delta[n-1]+3\delta[n-3]+9\delta[n-5],\;}
\]
\[
\{y_2[n]\}=\{1,\;\underset{\uparrow}{0},\;-5,\;0,\;3,\;0,\;9\},\qquad -1\le n\le 5 .
\]
(Term by term: $(-\delta[n])\circledast(-\delta[n+1])=\delta[n+1]$;
$(-\delta[n])\circledast(2\delta[n-1])=-2\delta[n-1]$;
$(3\delta[n-2])\circledast(-\delta[n+1])=-3\delta[n-1]$, and these two combine to $-5\delta[n-1]$;
$(-\delta[n])\circledast(3\delta[n-3])=-3\delta[n-3]$ and
$(3\delta[n-2])\circledast(2\delta[n-1])=6\delta[n-3]$, combining to $+3\delta[n-3]$;
$(3\delta[n-2])\circledast(3\delta[n-3])=9\delta[n-5]$.)

\textbf{(c)} $y_3[n]=x_1[n]\circledast h_2[n]$:
\[
\boxed{\;y_3[n]=2\delta[n+2]-6\delta[n]-2\delta[n-2]+6\delta[n-4],\;}
\]
\[
\{y_3[n]\}=\{2,\;0,\;\underset{\uparrow}{-6},\;0,\;-2,\;0,\;6\},\qquad -2\le n\le 4 .
\]

\textbf{(d)} $y_4[n]=x_2[n]\circledast h_1[n]$:
\[
\boxed{\;y_4[n]=-\delta[n+3]+3\delta[n+1]+1.5\delta[n]-3.5\delta[n-2]-3\delta[n-4],\;}
\]
\[
\{y_4[n]\}=\{-1,\;0,\;3,\;\underset{\uparrow}{1.5},\;0,\;-3.5,\;0,\;-3\},\qquad -3\le n\le 4 .
\]

Index-range check: in each case the support runs from the sum of the starting indices to the
sum of the ending indices, and the lengths are $3+6-1=8$, $3+5-1=7$, $3+5-1=7$, $3+6-1=8$
respectively. \checkmark

\Prob{2.48\quad Largest sample of convolutions of rectangular pulses}

The four sequences are rectangular pulses (with $N>M$):
\[
\begin{aligned}
x[n]&=\mu[n]-\mu[n-N] &&=1 \text{ for } 0\le n\le N-1 &&\text{(length }N),\\
h[n]&=\mu[n]-\mu[n-M] &&=1 \text{ for } 0\le n\le M-1 &&\text{(length }M),\\
g[n]&=\mu[n-N]-\mu[n-M] &&=-1 \text{ for } M\le n\le N-1 &&\text{(length }N-M),\\
w[n]&=\mu[n+M]-\mu[n-N] &&=1 \text{ for } -M\le n\le N-1 &&\text{(length }N+M).
\end{aligned}
\]
\Mis The page prints $g[n]=\mu[n-N]-\mu[n-N]\equiv0$ and repeats
$y_3=y_4=g\circledast g$. The evident intent (and the form used in the electronic edition) is
$g[n]=\mu[n-N]-\mu[n-M]$ and $y_4=w\circledast w$; that is what is solved here.

\textbf{The general rule.} The convolution of two rectangular pulses of amplitudes $a,b$ and
lengths $L_1,L_2$ starting at $n_1,n_2$ is a trapezoid: it rises linearly, holds a plateau of
value $ab\min(L_1,L_2)$, then falls. The plateau occupies
\[
n_1+n_2+\min(L_1,L_2)-1 \;\le\; n \;\le\; n_1+n_2+\max(L_1,L_2)-1 ,
\]
which degenerates to a single peak sample when $L_1=L_2$. No convolution sum is needed ---
only the overlap count.

\begin{center}
\renewcommand{\arraystretch}{1.6}
\begin{tabular}{@{}cllll@{}}
\toprule
& sequence & extreme value & location(s)\\
\midrule
(a) & $y_1=x\circledast x$ & $N$ & $n=N-1$ (single peak)\\
(b) & $y_2=h\circledast h$ & $M$ & $n=M-1$ (single peak)\\
(c) & $y_3=g\circledast g$ & $N-M$ & $n=M+N-1$ (single peak)\\
(d) & $y_4=w\circledast w$ & $N+M$ & $n=N-M-1$ (single peak)\\
(e) & $y_5=x\circledast h$ & $M$ & $M-1\le n\le N-1$ (plateau)\\
(f) & $y_6=x\circledast g$ & $-(N-M)$ & $N-1\le n\le N+M-1$ (plateau, \emph{negative})\\
(g) & $y_7=x\circledast w$ & $N$ & $N-M-1\le n\le N-1$ (plateau)\\
\bottomrule
\end{tabular}
\end{center}

Remarks on individual entries:
\begin{itemize}[leftmargin=1.4em,itemsep=2pt]
\item (a),(b),(d) are autoconvolutions of equal-length pulses, so they are triangles with a
single apex at the midpoint of the support: e.g.\ for $y_4$, the support is
$-2M\le n\le 2N-2$ and the apex is at $n=\tfrac{(-2M)+(2N-2)}{2}=N-M-1$.
\item (c): $g$ has amplitude $-1$, so $g\circledast g$ is \emph{positive}, with apex
$(-1)^2(N-M)=N-M$ at $n=M+(N-1)=M+N-1$.
\item (f): all products are $(+1)(-1)=-1$, so every nonzero sample of $y_6$ is negative. Its
largest \emph{positive} sample is therefore $0$; the extreme (most negative) sample is
$-(N-M)$, listed above.
\item (e),(g): unequal lengths give a flat top of height equal to the \emph{shorter} length.
\end{itemize}

\Prob{2.49\quad Symmetry of convolutions of conjugate-symmetric / antisymmetric parts}

Let $c[n]=(a\circledast b)[n]=\sum_k a[k]b[n-k]$. Then
\[
c^*[-n]=\sum_k a^*[k]\,b^*[-n-k]
\;\overset{m=-k}{=}\;\sum_m a^*[-m]\,b^*[-(n-m)] .
\]
So the symmetry type of $c$ is the \emph{product} of the symmetry types of $a$ and $b$, where
conjugate symmetric counts as $+1$ and conjugate antisymmetric as $-1$:
\[
a^*[-m]=\varepsilon_a\,a[m],\quad b^*[-(n-m)]=\varepsilon_b\,b[n-m]
\;\Longrightarrow\;
c^*[-n]=\varepsilon_a\varepsilon_b\sum_m a[m]b[n-m]=\varepsilon_a\varepsilon_b\,c[n].
\]

\textbf{(a)} $h_{\mathrm{cs}}\circledast g_{\mathrm{cs}}$: $\varepsilon=(+1)(+1)=+1$
$\Rightarrow$ \boxed{\text{conjugate symmetric}}

\textbf{(b)} $h_{\mathrm{ca}}\circledast g_{\mathrm{cs}}$: $\varepsilon=(-1)(+1)=-1$
$\Rightarrow$ \boxed{\text{conjugate antisymmetric}}

\textbf{(c)} $h_{\mathrm{ca}}\circledast g_{\mathrm{ca}}$: $\varepsilon=(-1)(-1)=+1$
$\Rightarrow$ \boxed{\text{conjugate symmetric}}

(Written out for (b): $c^*[-n]=\sum_m\bigl(-h_{\mathrm{ca}}[m]\bigr)\bigl(g_{\mathrm{cs}}[n-m]\bigr)=-c[n]$.)
